\documentclass{article}
\usepackage{polynom}
\usepackage{nicematrix,tikz}
\usetikzlibrary{fit}
\newcommand{\ff}{\mathtt{f}}

\pagestyle{empty}
\begin{document}
\begin{equation*}
\boldsymbol{\mathcal{F}}^{i} =
\begin{bNiceMatrix}[margin,nullify-dots]
    \ff_{1\,1} & \Cdots & \ff_{1\,\varpi} & \ff_{1\,(\varpi+1)} & \Cdots & \ff_{1\,\eta} \\
    \Vdots & \Ddots & \Vdots & \Vdots & & \Vdots \\
    \ff_{\varpi\,1} & \Cdots & \ff_{\varpi\,\varpi} & \ff_{\varpi\,(\varpi+1)} & \Cdots & \ff_{\varpi\,\eta} \\
    \noalign{\kern2mm}
    \ff_{(\varpi+1)\,1} & \Cdots & \ff_{(\varpi+1)\,\varpi} & \ff_{(\varpi+1)\,(\varpi+1)} & \Cdots & \ff_{(\varpi+1)\,\eta} \\
    \Vdots & & \Vdots & \Vdots & \Ddots & \Vdots \\
    \ff_{\eta\,1} & \Cdots & \ff_{\eta\,\varpi} & \ff_{\eta\,(\varpi+1)} & \Cdots & \ff_{\eta\,\eta}
\CodeAfter
  \SubMatrix[{1-1}{3-3}][slim,extra-height=-1mm,xshift=1mm]
  \OverBrace{1-1}{3-3}{\scriptstyle\mathcal{F}_N^i}[shorten,yshift=1mm]
\end{bNiceMatrix}
\end{equation*}
\end{document}